\documentclass[12pt]{article}
\usepackage[a4paper,margin=1in]{geometry}
\usepackage{amsmath,amssymb}
\usepackage{mathtools}
\usepackage{bm}

\title{Closed-Form Solution for Nonzero Hopping, Zero Interaction, and Nonzero Dissipation}
\author{}
\date{}

\begin{document}
\maketitle

\section*{Purpose}
This note derives the closed-form solution for a bosonic lattice with nonzero hopping amplitude $J$, zero onsite interaction $U$, and nonzero single-particle loss rate $\gamma$. The key point is that when $U=0$, the model remains quadratic, so the equations of motion for the first moments and the one-body correlation matrix close exactly.

\section*{Model}
Consider bosonic annihilation operators $\hat a_j$ on sites $j=1,\dots,M$. The Hamiltonian is
\[
\hat H
=
-J\sum_{\langle i,j\rangle}
\left(\hat a_i^\dagger \hat a_j + \hat a_j^\dagger \hat a_i\right),
\]
and the local loss operators are
\[
\hat L_j = \sqrt{\gamma}\,\hat a_j.
\]
The density matrix evolves according to the Lindblad master equation
\[
\frac{d\rho}{dt}
=
-i[\hat H,\rho]
+
\sum_{j=1}^M
\gamma
\left(
\hat a_j \rho \hat a_j^\dagger
-\frac12 \{\hat a_j^\dagger \hat a_j,\rho\}
\right).
\]

Since $U=0$, the Hamiltonian is quadratic. This is why the problem remains analytically solvable.

\section*{Quadratic Hamiltonian in Matrix Form}
Write the Hamiltonian as
\[
\hat H = \sum_{m,n} h_{mn}\,\hat a_m^\dagger \hat a_n,
\]
where $h$ is the single-particle hopping matrix. For nearest-neighbor hopping,
\[
h_{mn} =
\begin{cases}
-J, & m,n \text{ nearest neighbors},\\
0, & \text{otherwise}.
\end{cases}
\]

\section*{Equation of Motion for the First Moments}
For any operator $\hat O$, the Lindblad equation gives
\[
\frac{d}{dt}\langle \hat O\rangle
=
-i\langle[\hat O,\hat H]\rangle
+
\sum_j \gamma
\left\langle
\hat a_j^\dagger \hat O \hat a_j
-\frac12\{\hat a_j^\dagger \hat a_j,\hat O\}
\right\rangle.
\]

Take $\hat O=\hat a_\ell$. First compute the Hamiltonian commutator. Using
\[
[\hat a_\ell,\hat a_m^\dagger \hat a_n] = \delta_{\ell m}\hat a_n,
\]
we obtain
\[
[\hat a_\ell,\hat H]
=
\sum_{m,n} h_{mn}[\hat a_\ell,\hat a_m^\dagger \hat a_n]
=
\sum_n h_{\ell n}\hat a_n.
\]
Therefore, the coherent part contributes
\[
-i\langle[\hat a_\ell,\hat H]\rangle
=
-i\sum_n h_{\ell n}\langle \hat a_n\rangle.
\]

For the dissipative part, one finds
\[
\sum_j \gamma
\left\langle
\hat a_j^\dagger \hat a_\ell \hat a_j
-\frac12\{\hat a_j^\dagger \hat a_j,\hat a_\ell\}
\right\rangle
=
-\frac{\gamma}{2}\langle \hat a_\ell\rangle.
\]
Hence,
\[
\boxed{
\frac{d}{dt}\langle \hat a_\ell\rangle
=
-i\sum_n h_{\ell n}\langle \hat a_n\rangle
-\frac{\gamma}{2}\langle \hat a_\ell\rangle.
}
\]

\section*{Vector Form of the Solution}
Define the first-moment vector
\[
\bm{\alpha}(t)
=
\begin{pmatrix}
\langle \hat a_1(t)\rangle\\
\langle \hat a_2(t)\rangle\\
\vdots\\
\langle \hat a_M(t)\rangle
\end{pmatrix}.
\]
Then the equations of motion become
\[
\frac{d}{dt}\bm{\alpha}(t)
=
\left(-ih-\frac{\gamma}{2}I\right)\bm{\alpha}(t),
\]
where $I$ is the $M\times M$ identity matrix.

This is a linear ordinary differential equation with constant coefficients, so its solution is
\[
\boxed{
\bm{\alpha}(t)
=
e^{\left(-ih-\frac{\gamma}{2}I\right)t}
\bm{\alpha}(0).
}
\]
Because $I$ commutes with $h$, this can be factorized as
\[
\boxed{
\bm{\alpha}(t)
=
e^{-\gamma t/2}\,e^{-iht}\,\bm{\alpha}(0).
}
\]

This expression gives the exact closed-form evolution of the first-order coherence amplitudes.

\section*{Diagonalization of the Hopping Matrix}
If the hopping matrix is diagonalizable, write
\[
h = V \varepsilon V^{-1},
\]
where
\[
\varepsilon = \operatorname{diag}(\varepsilon_1,\dots,\varepsilon_M)
\]
contains the single-particle eigenvalues. Then
\[
e^{-iht} = V e^{-i\varepsilon t} V^{-1},
\]
so the solution becomes
\[
\boxed{
\bm{\alpha}(t)
=
e^{-\gamma t/2}
V\,\operatorname{diag}(e^{-i\varepsilon_1 t},\dots,e^{-i\varepsilon_M t})\,V^{-1}\bm{\alpha}(0).
}
\]

This makes the physical structure transparent:
\begin{itemize}
\item the hopping produces coherent oscillation and mode mixing through the phases $e^{-i\varepsilon_k t}$,
\item the dissipation contributes the universal damping factor $e^{-\gamma t/2}$.
\end{itemize}

\section*{Equation of Motion for the One-Body Correlation Matrix}
Now define the one-body correlation matrix
\[
C_{mn}(t)=\langle \hat a_m^\dagger \hat a_n\rangle.
\]
In matrix form,
\[
C(t) =
\begin{pmatrix}
\langle \hat a_1^\dagger \hat a_1\rangle & \cdots & \langle \hat a_1^\dagger \hat a_M\rangle\\
\vdots & \ddots & \vdots\\
\langle \hat a_M^\dagger \hat a_1\rangle & \cdots & \langle \hat a_M^\dagger \hat a_M\rangle
\end{pmatrix}.
\]

For a quadratic Hamiltonian and linear loss, the correlation matrix also obeys a closed equation:
\[
\boxed{
\frac{dC}{dt}
=
-i[h,C]-\gamma C.
}
\]
This has the exact solution
\[
\boxed{
C(t)
=
e^{\left(ih-\frac{\gamma}{2}I\right)t}
\,C(0)\,
e^{\left(-ih-\frac{\gamma}{2}I\right)t}.
}
\]
Again using the fact that $I$ commutes with $h$, this simplifies to
\[
\boxed{
C(t)
=
e^{-\gamma t}\,
e^{iht}\,
C(0)\,
e^{-iht}.
}
\]

This is the exact closed-form result for the one-body density matrix, which also contains the first-order coherence information.

\section*{Site Occupations}
The local particle number at site $j$ is
\[
\hat n_j = \hat a_j^\dagger \hat a_j,
\]
so its expectation value is obtained from the diagonal of the correlation matrix:
\[
\boxed{
\langle \hat n_j(t)\rangle = C_{jj}(t).
}
\]
Thus, once the matrix $C(t)$ is known, all site occupations follow directly.

\section*{Normal-Mode Picture}
Define normal-mode operators
\[
\hat b_k = \sum_j (V^{-1})_{kj}\hat a_j,
\]
so that the Hamiltonian becomes diagonal:
\[
\hat H = \sum_k \varepsilon_k \hat b_k^\dagger \hat b_k.
\]
Then each mode evolves independently. For the first moments,
\[
\boxed{
\langle \hat b_k(t)\rangle
=
e^{-\gamma t/2}e^{-i\varepsilon_k t}\langle \hat b_k(0)\rangle.
}
\]
For the mode correlations,
\[
\boxed{
\langle \hat b_k^\dagger(t)\hat b_q(t)\rangle
=
e^{-\gamma t}e^{i(\varepsilon_k-\varepsilon_q)t}
\langle \hat b_k^\dagger(0)\hat b_q(0)\rangle.
}
\]

This shows that in the normal-mode basis the exact solution consists simply of coherent oscillation plus exponential damping.

\section*{Interpretation}
When $J\neq 0$, $U=0$, and $\gamma\neq 0$:
\begin{itemize}
\item hopping transports amplitude and coherence between sites,
\item the absence of interaction means there is no nonlinear many-body mixing,
\item local dissipation damps all amplitudes exponentially,
\item and the full problem reduces to linear algebra on the single-particle hopping matrix.
\end{itemize}

This is why a closed-form solution exists in this case, whereas for $U\neq 0$ such a simple exact solution is generally no longer available.

\section*{Conclusion}
For nonzero hopping, zero interaction, and nonzero dissipation, the bosonic lattice remains exactly solvable because the Hamiltonian is quadratic. The exact solutions are
\[
\bm{\alpha}(t)
=
e^{-\gamma t/2}e^{-iht}\bm{\alpha}(0)
\]
for the first moments, and
\[
C(t)
=
e^{-\gamma t}e^{iht}C(0)e^{-iht}
\]
for the one-body correlation matrix. These formulas provide the closed-form evolution of first-order coherence and local occupations throughout the lattice.

\end{document}
